\section{From kernel vectors to possible columns}\label{sec:reductions}

We reduce Theorem~\ref{thm:main} to one estimate on the number of
columns that can complete a nonsingular bit matrix to a singular one.
We first treat kernel vectors whose coordinates away from a largest
one, after normalization, are close to a vector supported on a small
set. For every remaining kernel vector, we will delete
a coordinate that is large enough to control the inverse matrix and
whose original row has the expected proportion of positive signs.
The latter condition supplies the entropy needed for the biased law.

Negating the entire matrix replaces \(p\) by \(1-p\) and preserves
singularity. We therefore assume throughout the proof that \(0<p<1/2\),
and set
\begin{equation}\label{eq:rate-constants}
 b=1-p,\qquad Q=p^2+b^2,\qquad
 a=-\log_2 b,\qquad d_*=-\log_2 Q,\qquad
 h_*=-p\log_2 p-b\log_2 b.
\end{equation}
Two strict inequalities will be used:
\begin{equation}\label{eq:rate-margins}
 2a>d_*,\qquad h_*>d_*.
\end{equation}
The first follows from \(b^2<Q\). For the second, strict concavity of
the logarithm gives
\(p\log p+b\log b<\log(p^2+b^2)\).
The use of column counts to exploit this entropy gap has precedents
in independent-entry models
\cite[Section~2.1]{LitvakTikhomirov};
the comparison itself is explicit in
Jain, Sah, and Sawhney \cite[Section~6]{JainSahSawhneyDiscrete}.
Here it will be applied to an original row selected from the kernel
coordinates of a symmetric matrix.
Constants may depend on \(p\) and on any other parameters fixed before
the dimension tends to infinity. In \(O_p(\cdot)\), the implicit
constant depends only on \(p\). The logarithm \(\log\) is natural,
unless a different base is displayed, and \([m]=\{1,\ldots,m\}\).

\subsection{Corank one and tails close to sparse vectors}

The diagonal-flip reduction of
\cite[Lemma~2.4]{SymmetricRate} remains valid under the biased law,
with a subexponential change in probability.

\begin{lemma}\label{lem:corank-reduction}
For the symmetric sign matrix \(A=A_n^{(p)}\),
\[
 \Prob(\det A=0)
 \le 2^{o(n)}\Prob(\rank A=n-1)+2^{-\Omega_p(n^{6/5})}.
\]
\end{lemma}

\begin{proof}
A symmetric matrix of rank \(r\) has a nonsingular principal
submatrix of size \(r\). Indeed, write it as \(XSX^{\mathsf T}\),
where \(X\) has \(r\) independent columns and \(S\) is nonsingular,
and choose \(r\) independent rows of \(X\). The principal submatrix
on those rows is a product of three nonsingular square matrices.
For a fixed candidate principal submatrix, reveal it and its cross
block. If the full matrix has the same rank, the Schur complement
forces all entries in the remaining principal block. Each forced
entry has conditional probability at most \(b\). Consequently,
with \(t=\lceil n^{3/5}\rceil\),
\[
 \Prob(n-\rank A\ge t)
 \le \sum_{k=t}^n\binom nk b^{k(k+1)/2}
 \le 2^{-\Omega_p(n^{6/5})}.
\]

Now let \(K=\ker A\) have dimension \(1\le k<t\). Choose a set
\(I\) of \(k-1\) coordinates whose restrictions to \(K\) are
independent linear functionals, and flip the diagonal signs on \(I\).
Write \(D\) for the diagonal matrix of these changes. If
\((A+D)y=0\), then
\[
 x^{\mathsf T}Dy=0\qquad(x\in K).
\]
The map \(x\mapsto x_I\) maps \(K\) onto \(\R^I\), and each entry
of \(D\) on \(I\) is nonzero. Hence \(y_I=0\), and then \(Ay=0\).
Conversely, every vector in \(K\) vanishing on \(I\) is in
\(\ker(A+D)\). The new kernel is exactly that one-dimensional
subspace.

Fix a choice of \(I\) for each such matrix. The probability of a
matrix divided by that of its image is at most \((b/p)^t\).
Each image has at most
\[
 \sum_{j\le t}\binom nj=2^{O(n^{3/5}\log n)}
\]
preimages, since specifying the flipped coordinates recovers the
original matrix. Both factors are \(2^{o(n)}\), proving the lemma.
\end{proof}

Given a nonzero kernel vector \(v\), call an index \(a_0\) with
\(|v_{a_0}|=\max_i|v_i|\) an \emph{anchor}. The vector off the
anchor cannot vanish, since no column of \(A\) has a zero entry.
We normalize it by
\[
 \lVert v_{[n]\setminus\{a_0\}}\rVert_2=1
\]
and call this normalized off-anchor vector the \emph{tail}.
Every row equation \(Av=0\) gives
\begin{equation}\label{eq:anchor-size}
 |v_{a_0}|
 \le \sum_{j\ne a_0}|v_j|
 \le \sqrt{n-1}.
\end{equation}

\begin{lemma}[Tails close to sparse vectors]\label{lem:compressible}
For every \(\epsilon>0\), there are fixed
\(\lambda,\delta\in(0,1/4)\) such that, for all sufficiently large
\(n\), the following event has probability at most
\(2^{(-d_*+\epsilon)n}\): the matrix \(A_n^{(p)}\) has a nonzero
kernel vector whose normalized tail, for some anchor, lies within
Euclidean distance \(\delta\) of a vector supported on at most
\(\lambda n\) coordinates.
\end{lemma}

\begin{proof}
Put \(\mu=2p-1\) and
\(W=A-\mu\mathbf1\mathbf1^{\mathsf T}\). For a sufficiently large
fixed \(L\),
\begin{equation}\label{eq:centered-norm}
 \Prob(\lVert W\rVert_{\mathrm{op}}>L\sqrt n)\le 2^{-4n}.
\end{equation}
To see this, for a fixed unit vector \(x\), write \(x^{\mathsf T}Wx\)
as a sum over the independent upper-triangular entries. Its
coefficients have squared sum
\[
 \sum_i x_i^4+4\sum_{i<j}x_i^2x_j^2\le2.
\]
Hoeffding's inequality \cite[Theorem~2]{Hoeffding} therefore gives
a bound \(2\exp(-cL^2n)\) at a constant multiple of \(L\sqrt n\).
A \(1/4\)-net of the unit sphere has at most \(9^n\) points, and
the usual quadratic-form net inequality bounds the operator norm
by twice the maximum over this net. Taking \(L\) large proves
\eqref{eq:centered-norm}.

Fix an anchor \(a_0\) and a tail support \(T\), with
\(j=|T|\le\lambda n\). If a tail is within \(\delta\) of a vector
on \(T\), its truncation to \(T\) has error at most \(\delta\).
Approximate this truncation by a \(\delta\)-net of the unit ball
in \(\R^T\), and approximate the anchor coefficient on
\([-\sqrt n,\sqrt n]\) with mesh at most \(\delta\).
The resulting vector \(x\), supported on
\(S=T\cup\{a_0\}\), satisfies
\[
 \lVert v-x\rVert_2\le3\delta,\qquad
 \lVert x_T\rVert_2\ge1-2\delta,\qquad
 |x_{a_0}|\ge\max_{i\in T}|x_i|-2\delta.
\]
It suffices to retain net points with these last two properties.
Also approximate \(\sum_i v_i\) by a number \(u\) with error at
most \(\delta\). By \eqref{eq:anchor-size}, this sum lies in
\([-2\sqrt n,2\sqrt n]\).
For each \(T\), the number of pairs \((x,u)\) is at most
\((C/\delta)^j\) times a polynomial in \(n\), with fixed parameters.
The separate scalar approximation prevents the mean matrix from
amplifying the error in the coordinate sum.

On the event \(\lVert W\rVert_{\mathrm{op}}\le L\sqrt n\),
the equation \(Av=0\)
implies that one of these pairs satisfies
\begin{equation}\label{eq:net-residual}
 \lVert Wx+\mu u\mathbf1\rVert_2
 \le C_1L\delta\sqrt n.
\end{equation}
We next show that every fixed pair has small probability of
satisfying this inequality. There are fixed \(w>0\) and
\(\delta_0>0\), depending on \(p\), such that, whenever
\(\delta<\delta_0\), each row outside \(S\) satisfies
\begin{equation}\label{eq:two-sign-concentration}
 \Prob\bigl(|(Wx)_i+\mu u|\le w\bigr)\le Q.
\end{equation}
All row shifts here may be arbitrary.

Choose a small fixed coefficient threshold \(u_0>0\). If some
\(|x_j|\ge u_0\), \(j\in T\), then
\(|x_{a_0}|\ge u_0/2\) when \(\delta\) is small enough.
Condition on all signs in this row except those at \(j,a_0\).
For \(w<u_0/4\), the interval of length \(2w\) cannot contain
two outcomes that differ in only one of these two signs.
The allowed vertices of the two-bit square thus consist of at
most one vertex or two opposite vertices. Their total probability
is at most
\[
 \max\{b^2,p^2,2pb,p^2+b^2\}=Q.
\]
If instead every tail coefficient has absolute value below \(u_0\),
condition on the anchor sign and use the tail alone.
Its centered sum has variance
\[
 \sigma_T^2=4pb\sum_{j\in T}x_j^2,
\]
bounded below by a positive constant depending on \(p\).
The sum of its third absolute centered moments is at most
\(C_pu_0\sum_{j\in T}x_j^2\). The Berry--Esseen inequality
\cite{Berry,Esseen}, in its unequal-summand form
\cite[Theorem, p.~586]{Tyurin}, bounds the uniform
distribution-function error by an absolute constant times the
sum of these third moments divided by \(\sigma_T^3\).
The error from the centered normal law with variance
\(\sigma_T^2\) is therefore at most \(C_pu_0\). The probability of any
interval of length \(2w\) is therefore at most \(C_p(w+u_0)\).
Choose \(u_0\), then \(w\), small enough to make this quantity
less than \(Q\), and finally decrease \(\delta_0\). This proves
\eqref{eq:two-sign-concentration} in both cases.

For a fixed pair \((x,u)\), the rows outside \(S\) depend on
disjoint sets of upper-triangular entries, so these tests are
independent. Choose a small fixed \(\gamma>0\), and then
\(\delta<\delta_0\) so that
\(C_1^2L^2\delta^2/w^2<\gamma\).
Inequality~\eqref{eq:net-residual} allows at most \(\gamma n\)
rows to fail the interval test. Thus its probability, for fixed
data, is at most
\[
 \left(\sum_{i\le\gamma n}\binom ni\right)
 Q^{\,n-j-1-\gamma n}.
\]
The norm event was used only to obtain \eqref{eq:net-residual};
the independent tests in this bound are not conditioned on that
event.

For \(t\in(0,1)\), let
\(H_2(t)=-t\log_2t-(1-t)\log_2(1-t)\).
Union over anchors, supports, and net points gives a base-two
logarithm at most
\[
 -d_*n+
 \bigl[d_*(\lambda+\gamma)+H_2(\gamma)+H_2(\lambda)
       +\lambda\log_2(C/\delta)\bigr]n+o(n).
\]
Choose \(\gamma\) sufficiently small, then \(\delta\) as above,
then \(\lambda\) sufficiently small, to make the bracket less
than \(\epsilon/2\). Adding \eqref{eq:centered-norm} and taking
\(n\) large proves the lemma.
\end{proof}

\subsection{The conditioned bit model and the column estimate}

For tails outside Lemma~\ref{lem:compressible}, we will use
Schur complementation to identify the possible deleted columns.
The sign-to-bit congruence is the one in
\cite[Lemma~2.1]{SymmetricRate}; conditioning must now retain its
effect on the entry probabilities.

\begin{lemma}\label{lem:anchor-congruence}
Fix an index \(a_0\) and condition on its entire row in
\(A_n^{(p)}\). There is a symmetric bit matrix \(M\) of size
\(N=n-1\) with the following properties:
\begin{enumerate}[label=\textup{(\roman*)}]
\item the entries on and above its diagonal are independent,
and each has success probability \(p\) or \(b\);
\item \(\rank(\E M)\le2\), and the same holds for every
principal minor;
\item \(\dim\ker M=\dim\ker A\); whenever \(a_0\) is an anchor
of a kernel vector of \(A\), its normalized tail, with some
coordinate signs changed, is a unit vector in \(\ker M\).
\end{enumerate}
\end{lemma}

\begin{proof}
Relabel the fixed index as the first coordinate, put \(\mu=2p-1\),
\(t=A_{11}\)
and \(u_i=A_{1i}\) for \(i>1\), and set
\[
 D=\diag(1,tu_2,\ldots,tu_n).
\]
The matrix \(tDAD\) has first row and column equal to one.
Its lower principal block is the sign matrix
\(C_{ij}=tu_i u_j A_{ij}\), \(i,j>1\). Hence its Schur complement
at the first entry is
\[
 C-\mathbf1\mathbf1^{\mathsf T}=-2M,\qquad
 M_{ij}=\frac{1-tu_i u_j A_{ij}}2.
\]
After the anchor row is fixed, the remaining original signs are
independent. The displayed formula proves \textup{(i)} and gives
\[
 \E M=\frac12\mathbf1\mathbf1^{\mathsf T}
       -\frac{\mu t}{2}uu^{\mathsf T}.
\]
This includes the diagonal, since \(u_i^2=1\), and proves
\textup{(ii)}. Multiplication by \(t\), diagonal congruence,
and taking the Schur complement of the nonzero first entry
preserve nullity. More explicitly, \(Dv\) is a kernel vector
of \(tDAD\), and its lower coordinates are in \(\ker M\).
They differ from the original tail only by coordinate signs,
which proves \textup{(iii)}.
\end{proof}

For the rest of the paper, let \(N\) be a reference dimension and put
\begin{equation}\label{eq:reference-range}
 \ell=\log\log N,\qquad
 \mathcal I_N=\{m\in\Z:N-N^{7/10}\le m\le N\}.
\end{equation}
By the \emph{bit model} we mean a symmetric matrix \(B\) of size
\(m\in\mathcal I_N\), with independent entries on and above the
diagonal, whose fixed success probabilities belong to
\(\{p,b\}\), and with \(\rank(\E B)\le2\).
All estimates for this model are uniform over these sizes and
arrays of probabilities. A property \emph{holds typically} if
its failure probability is \(2^{-\omega(N)}\) uniformly in this
sense. For a property stated only for nonsingular matrices,
failure means nonsingularity together with failure of the
conclusion. Any fixed parameters in the property are chosen
before \(N\) tends to infinity.

For \(L_0\ge1\) fixed and a nonsingular matrix \(B\) of size \(m\),
define
\begin{equation}\label{eq:admissible-columns}
 \mathcal S_{L_0}(B)=
 \left\{z\in\bits^m:
 z^{\mathsf T}B^{-1}z\in\{0,1\},\
 \lVert B^{-1}z\rVert_2\le L_0\sqrt N\right\}.
\end{equation}
The quadratic condition is exactly the Schur-complement condition
for a bit diagonal to complete \(B\) to a singular matrix.
The norm bound will follow from the choice of deleted coordinate.
The only deferred input in the proof of the main theorem is the
following estimate.

\begin{proposition}[Column estimate]\label{prop:column-bound}
For every fixed \(L_0\ge1\) and every sufficiently small fixed
\(e_0>0\), a matrix \(B\) in the bit model of size \(m=N-1\)
satisfies, for all sufficiently large \(N\),
\[
 \Prob\left(
 B\text{ is nonsingular and }
 |\mathcal S_{L_0}(B)|>2^{e_0N}
 \right)
 \le 2^{-N(2a-10e_0)}.
\]
\end{proposition}

Proposition~\ref{prop:column-bound} will be proved at the end of
Section~\ref{sec:paths}. Its cardinality bound is useful here
because a row with approximately \(pn\) positive original signs
has point probability close to \(2^{-h_*n}\). We next ensure
that such a row can be chosen among the large kernel coordinates.

\begin{lemma}[Row counts]\label{lem:row-counts}
For every fixed \(\rho,w>0\), with probability
\(1-\exp(-\Omega_{\rho,w}(n^2))\), fewer than \(\rho n\) rows
of \(A_n^{(p)}\) have a number of positive entries outside
\([pn-wn,pn+wn]\).
\end{lemma}

\begin{proof}
If at least \(\rho n\) rows fail, at least half fail in the same
direction. Choose a set \(I\) of \(k=\lceil\rho n/3\rceil\)
such rows, for all sufficiently large \(n\); it suffices to
consider \(\rho<1\).
The sum of their positive-entry counts is a sum of independent
upper-triangular Bernoulli variables with coefficients one or
two. Its mean is \(pkn\), and its coefficients have squared
sum at most \(2kn\). The common-direction deviation is at
least \(kwn\). Hoeffding's inequality
\cite[Theorem~2]{Hoeffding} bounds this probability by
 \(2\exp(-kw^2n)\).
There are at most \(2^n\) sets \(I\) and two directions, giving
the assertion.
\end{proof}

\begin{proof}[Proof of Theorem~\ref{thm:main}, assuming
Proposition~\ref{prop:column-bound}]
Fix \(\epsilon>0\), and choose \(\lambda,\delta\) from
Lemma~\ref{lem:compressible}. By
Lemma~\ref{lem:corank-reduction}, it suffices to bound the
corank-one event; the matrices discarded in the corank reduction
are already negligible. Discard also the event in
Lemma~\ref{lem:compressible}.

Take any remaining normalized tail and choose a fixed
\(c\in(0,\delta)\). There are at least \(\lambda n/2\)
off-anchor coordinates with absolute value at least
\(c/\sqrt n\). Otherwise truncation to these coordinates
would have support at most \(\lambda n\) and discarded norm
at most \(c\), contrary to the choice of the tail.
For a fixed tolerance \(w_1>0\), Lemma~\ref{lem:row-counts}
with \(\rho=\lambda/4\) shows, apart from probability
\(\exp(-\Omega_{\lambda,w_1}(n^2))\), that some such coordinate
has a row count within \(w_1n\) of \(pn\).
Choose it as the deletion index.

Fix the anchor and deletion indices and condition on the anchor
row. This does not condition on their having been chosen by a
kernel vector; we will take a union bound over their indices.
Apply Lemma~\ref{lem:anchor-congruence}, and denote the transformed
unit tail by \(v\in\ker M\), where \(M\) has size \(N=n-1\)
and corank one. Its deletion coordinate has
\(|v_i|\ge c/\sqrt n\).
The adjugate of a symmetric corank-one matrix is a nonzero
scalar multiple of \(vv^{\mathsf T}\). Thus deleting coordinate
\(i\) leaves a nonsingular principal minor \(B\) of size \(N-1\).
Writing
\[
 M=\begin{pmatrix}B&z\\z^{\mathsf T}&d\end{pmatrix},
 \qquad d\in\{0,1\},
\]
the kernel equation and Schur complement give
\[
 B^{-1}z=-\frac{v_{[N]\setminus\{i\}}}{v_i},
 \qquad z^{\mathsf T}B^{-1}z=d.
\]
Consequently \(z\in\mathcal S_{L_0}(B)\) for a fixed
\(L_0\ge\max(1,2/c)\).

It remains to bound the probability of these possible columns
under the biased law. Conditional on the anchor row, \(B\)
and the entries of \(z\) are independent. A specified \(z\)
specifies the \(n-2=N-1\) original off-diagonal signs of the
deleted column outside the anchor. If \(r\) of those signs
are positive, its point probability is exactly
\[
 p^r b^{N-1-r}.
\]
The full original row has only two additional entries, its
diagonal and its anchor entry. Compatibility with the selected
row-count condition therefore requires
\(|r-pn|\le w_1n+2\).
For every such \(z\),
\begin{equation}\label{eq:column-entropy}
 p^r b^{N-1-r}
 \le 2^{-N(h_*-C_pw_1-o(1))}.
\end{equation}
Indeed its negative base-two logarithm is
\[
 (N-1)h_*+
 \bigl(r-p(N-1)\bigr)\log_2(b/p).
\]
The count here concerns positive \emph{original signs};
the success probabilities of the transformed bits need not
be equal. We have not charged the deleted diagonal: dropping
its requirement only enlarges the event.

On the event that the nonsingular \(B\) satisfies the cardinality
conclusion of Proposition~\ref{prop:column-bound}, independence
and \eqref{eq:column-entropy} bound the conditional probability
by
\[
 2^{-N(h_*-e_0-C_pw_1-o(1))}.
\]
The exceptional nonsingular minors have conditional probability
at most \(2^{-N(2a-10e_0)}\). Both bounds are uniform in the
anchor row, so averaging over that row introduces no additional
factor. The choices of anchor and deletion indices cost at
most \(n^2\).

After \(L_0\) has been fixed, choose \(e_0,w_1>0\) sufficiently
small that
\[
 2a-10e_0>d_*,
 \qquad h_*-e_0-C_pw_1>d_*,
\]
using \eqref{eq:rate-margins}. Together with the discarded
events and the subexponential corank reduction, this proves
\[
 \limsup_{n\to\infty}\frac1n\log_2
       \Prob(\det A_n^{(p)}=0)\le-d_*+\epsilon.
\]
Letting \(\epsilon\downarrow0\) gives the required upper bound.

For the lower bound, require the first two rows to be equal.
The entries in their first two coordinates must satisfy
\(A_{11}=A_{12}=A_{22}\), at cost \(p^3+b^3\).
For each \(j\ge3\), the independent pair
\((A_{1j},A_{2j})\) agrees with probability \(Q\).
These requirements involve disjoint upper-triangular entries,
and hence
\[
 \Prob(\det A_n^{(p)}=0)\ge(p^3+b^3)Q^{n-2}.
\]
This gives the matching lower bound. Negation handles \(p>1/2\).
\end{proof}
